% ============================================
% LUC-07 (ZH)
% 中文翻译 — SPEC-935-R356
% ============================================
\section*{\textbf{LUC-07} --- 奥利维拉的踪迹}
\markboth{LUC-07 --- 奥利维拉的踪迹}{}

\noindent \textit{沿反向循环追踪}

\noindent \textbf{2026年4月 --- 贝洛奥里藏特 / 基谢拉莫宾}

\vspace{0.5em}\noindent\rule{\textwidth}{0.5pt}\vspace{0.5em}

\noindent 读完奥利维拉的信后，卢西亚做了任何法医精神病学家都会做的事：追踪踪迹。

\noindent 埃托尔·奥利维拉医生于1979年6月13日失踪。三周后，他的车在塞阿拉州基谢拉莫宾的一个停车场被发现。汽油耗尽了。车门没锁。车里有：一个装着衣服的背包，一本写了32页的笔记本（没人读过），还有一张6月17日去巴尔巴塞纳的巴士票——正是他到达的那天。

\noindent 他回去了。

\noindent 卢西亚在公共档案馆花了一下午翻箱子。她找到了奥利维拉的笔记本——被遗忘的，落满灰尘的，标签上写着\textquotedblleft{}H. 奥利维拉医生遗物 --- 1979年\textquotedblright{}。

\noindent 她40分钟读完了那32页。

\noindent 是临床记录，但有一个不同之处：笔迹在页面间变化。开始是坚定的、医生的笔迹。结束时颤抖、歪斜、\textit{幼稚}。

\noindent 最后一句话是：

\noindent \textit{\textquotedblleft{}我回到了巴尔巴塞纳。循环不让我走。我把日记还给档案馆。希望下一位读者能比我理解得更好。\textquotedblright{}}

\noindent \textbf{下一位读者。}

\noindent 她。

\noindent 她就是下一位。

\noindent 队列已经轮到她。

\vspace{0.5em}\noindent\rule{\textwidth}{0.5pt}\vspace{0.5em}

\noindent 卢西亚在档案馆过了一夜。

\noindent 她读了奥利维拉的日记。读了会谈报告。读了奥利维拉留给她的信——1979年写成，2026年送达。

\noindent 然后，第二次读了若阿金的日记。

\noindent 只不过这一次，她换了一双眼睛去读。

\noindent 不是作为精神病学家。
\noindent 而是作为\textit{继承人}。

\noindent 每个词都钻进来。每句话都回响。她感到甜味从纸页上升起。右眼开始灼烧。

\noindent \textit{\textquotedblleft{}循环没有回头路\textquotedblright{}}，她体内一个声音说。

\noindent 那不是她的声音。

\noindent 她合上日记。

\noindent 放进背包。

\noindent 然后，第一次，她确信自己停不下来了。

\vspace{0.5em}
\noindent\textit{\textbf{↪ 路线:} \rota{LUC-08} \rotasep \rota{DOC-11}}
